% Modelo de gabarito. Copie para materias/<matéria>/estudo/simulados/.fontes/<AAAA-MM-DD>-<prova>-gabarito.tex,
% ajuste os dados do cabeçalho e escreva as resoluções no lugar dos marcadores.
% Tudo o que está entre \marcador{...} e a figura de mostruário são exemplos de uso: não aproveite o conteúdo.
\documentclass[11pt]{article}
\usepackage[a4paper,left=2.2cm,right=2.2cm,top=3.5cm,bottom=2.4cm,headheight=50pt,headsep=12pt]{geometry}
\usepackage[T1]{fontenc}
\usepackage[brazil]{babel}
\usepackage[sc]{mathpazo}
\linespread{1.05}
\usepackage{amsmath,amssymb}
\usepackage{enumitem}
\usepackage{xcolor}
\usepackage{graphicx}
\usepackage{tabularx,booktabs}
\usepackage[breakable]{tcolorbox}
\usepackage{tikz}
\usetikzlibrary{arrows.meta,patterns,calc}
\usepackage{lastpage}
\usepackage{fancyhdr}
\usepackage[hidelinks,pdftitle={Simulado - resolução}]{hyperref}

% ---------------------------------------------------------------- dados do simulado
\newcommand{\materia}{PME 3100 -- MECÂNICA I}
\newcommand{\prova}{Prova 2}
\newcommand{\ano}{2026}

% ---------------------------------------------------------------- cabeçalho da Poli
\newcommand{\brasao}{\IfFileExists{../../../../../modelos/poli.png}%
  {\includegraphics[height=1.55cm]{../../../../../modelos/poli.png}}{}}
\pagestyle{fancy}\fancyhf{}
\fancyhead[L]{\brasao}
\fancyhead[C]{\large ESCOLA POLITÉCNICA DA UNIVERSIDADE DE SÃO PAULO\\[7pt] \materia\\[-3pt]}
\renewcommand{\headrulewidth}{0.6pt}
\renewcommand{\footrulewidth}{0.4pt}
\cfoot{\thepage}
\setlength{\parindent}{0pt}
\setlength{\parskip}{4pt}

% ---------------------------------------------------------------- notação e resolução
\newcommand{\vi}{\vec{\imath}}
\newcommand{\vj}{\vec{\jmath}}
\newcommand{\vk}{\vec{k}}
\newcommand{\marcador}[1]{\ifmmode\text{\color{gray}$\langle$\textit{#1}$\rangle$}\else\leavevmode{\color{gray}$\langle$\textit{#1}$\rangle$}\fi}
\definecolor{pontos}{RGB}{190,30,30}
\newcommand{\ok}{\;\checkmark}
% \resolucao{n}{valor}: título centrado de cada questão
\newcommand{\resolucao}[2]{\par\bigskip\begin{center}\underline{\textbf{Resolução -- Questão #1 (#2 pontos)}}\end{center}}
% \passo{a)}{Título}: abre o item; o texto segue na mesma linha
\newcommand{\passo}[2]{\par\smallskip\noindent\textbf{#1 #2.}\ }
% \pontos{0,5}: pontuação do passo, em vermelho na margem direita
\newcommand{\pontos}[1]{\hfill{\color{pontos}(#1)}}
% \resultado{<matemática>}{0,5}: resposta final em caixa, com a pontuação na margem
\newcommand{\resultado}[2]{\par\smallskip\noindent\fbox{$\displaystyle #1$}\pontos{#2}\par\smallskip}
% \res{<matemática>}: resposta em caixa dentro do texto
\newcommand{\res}[1]{\boxed{#1}}
% Observações do corretor: \begin{armadilhas} \item ... \end{armadilhas}
\newtcolorbox{armadilhasbox}{breakable, sharp corners, boxrule=0.5pt, colback=white, colframe=black!75,
  left=5pt, right=5pt, top=4pt, bottom=4pt}
\newenvironment{armadilhas}
  {\begin{armadilhasbox}\textbf{Observações do corretor (armadilhas desta questão).}
   \begin{enumerate}[label=\textbf{\arabic*.}, leftmargin=*, nosep]}
  {\end{enumerate}\end{armadilhasbox}}
% Tabela de critérios, lida pelo /corrigir: \criterios{ a) & ... \\ b) & ... \\ }
\newcommand{\criterios}[1]{\par\smallskip{\small
  \begin{tabularx}{\textwidth}{@{}lX@{}}\toprule \textbf{Item} & \textbf{Critério de correção} \\\midrule
  #1 \bottomrule\end{tabularx}}\par}
\newcommand{\inspirada}[1]{\par{\small\itshape Inspirada em: #1}\par}

% ---------------------------------------------------------------- kit TikZ (o mesmo do simulado)
% Na resolução, redesenhe a figura da questão e marque a construção com construcao e vetorres.
\tikzset{
  >=Stealth,
  every picture/.append style={line cap=round, line join=round, font=\small},
  solo/.style={pattern=north east lines, draw=none},
  corpo/.style={thick, fill=gray!12},
  barra/.style={line width=2.4pt},
  ponto/.style={circle, fill, inner sep=1.1pt},
  articulacao/.style={circle, draw, fill=white, inner sep=1.5pt},
  vetor/.style={->, thick},
  cota/.style={<->, thin},
  auxiliar/.style={densely dashed, gray},
  construcao/.style={densely dashed, pontos},
  vetorres/.style={->, thick, pontos},
  cavaleira/.style={x={(-0.45cm,-0.38cm)}, y={(1cm,0cm)}, z={(0cm,1cm)}},
  pics/eixos3d/.style={code={
    \draw[->] (0,0,0) -- (#1,0,0) node[pos=1.12] {$x$};
    \draw[->] (0,0,0) -- (0,#1,0) node[pos=1.1] {$y$};
    \draw[->] (0,0,0) -- (0,0,#1) node[pos=1.1] {$z$};}},
  pics/eixos3d/.default=1,
}
\newcommand{\piso}[2]{\fill[solo] ($#1+(0,-0.18)$) rectangle #2; \draw[thick] #1 -- #2;}
\newcommand{\parede}[3]{\fill[solo] #1 rectangle ($#2+(#3,0)$); \draw[thick] #1 -- #2;}
\newcommand{\giro}[5][]{\draw[->,#1] #2 ++(#3:#5) arc (#3:#4:#5)}

\begin{document}

\begin{center}
  {\bfseries \materia{} -- \ano{} -- \prova{} (Simulado) -- RESOLUÇÃO}\\[2pt]
  {\itshape Cada valor em vermelho é a pontuação do passo.
  \marcador{critério de desconto da prova real, tirado do mapa}}
\end{center}

% ---------------------------------------------------------------- questão 1
\resolucao{1}{3,0}

\marcador{Dados de partida: referenciais, vetores conhecidos, posições no instante.}

\begin{minipage}[t]{0.56\textwidth}
\passo{a)}{\marcador{Título do item}} \marcador{Raciocínio em uma ou duas frases, com a fórmula usada.}
\[ \vec v_B = \vec v_A + \vec\omega \wedge (B - A) \]
\resultado{\vec v_B = \marcador{resposta}}{0,5}
\end{minipage}\hfill
\begin{minipage}[t]{0.40\textwidth}\centering\vspace{0pt}
\begin{tikzpicture}
  % Mostruário da construção gráfica. Troque pela figura da questão com a construção marcada.
  \coordinate (A) at (0,0); \coordinate (B) at (2.2,1.2); \coordinate (I) at (0,2.2);
  \draw[barra] (A) -- (B);
  \node[articulacao, label=below:$A$] at (A) {}; \node[articulacao, label=right:$B$] at (B) {};
  \draw[vetorres] (A) -- ++(0.8,0) node[below] {$\vec v_A$};
  \draw[vetorres] (B) -- ++(0,0.7) node[right] {$\vec v_B$};
  \draw[construcao] (A) -- (I) -- (B);
  \node[ponto, pontos, label={[pontos]above:CIR}] at (I) {};
\end{tikzpicture}
\end{minipage}

\passo{b)}{\marcador{Título do item}} \marcador{Conta.}
\resultado{\marcador{resposta}}{1,0}
Conferência por outro caminho: \marcador{propriedade das projeções, outro polo, outro método}\ok

\begin{armadilhas}
  \item \marcador{O erro mais comum neste item, o que ele dá e por que está errado.}
  \item \marcador{O ``zero'' que é mesmo zero, ou o que parece zero e não é.}
\end{armadilhas}

\criterios{
  a) & \marcador{parcela} 0,25; \marcador{parcela} 0,25 \\
  b) & \marcador{parcela} 0,5; \marcador{parcela} 0,5. \marcador{Crédito parcial: vale 0,25 se...} \\
}
\inspirada{\marcador{prova e questão de origem, citadas do mapa, e o que mudou aqui}.}

% ---------------------------------------------------------------- mapa do simulado
\clearpage
\begin{center}\underline{\textbf{Mapa do simulado -- o que cada item treina}}\end{center}

{\small
\begin{tabularx}{\textwidth}{@{}lXp{4.6cm}@{}}
\toprule
\textbf{Item} & \textbf{Conteúdo / habilidade} & \textbf{Onde caiu antes} \\
\midrule
1a & \marcador{habilidade treinada} & \marcador{prova e questão, do mapa} \\
1b & \marcador{habilidade treinada} & \marcador{``nunca caiu'', se for o caso} \\
\bottomrule
\end{tabularx}}

\textbf{Como usar.} \marcador{Faça com cronômetro na duração da prova real, corrija pelos valores em vermelho e
anote, em cada item perdido, a observação do corretor que explica o erro. Diga quais itens refazer na véspera.}

\textbf{Autoavaliação.} \marcador{Faixas de nota com o que fazer em cada uma, citando as provas antigas a
refazer.}

\end{document}
